\subsubsection{Management Control}

Challenges, both expected and unexpected, are always present in any project. Even with a precisely estimated timeline and budget, it is very common for a project to fall short of the initial plan. 

At the end of the project, the real execution cost is calculated using the same methodology as the initial budget estimation. Based on these results, the project budget deviation is determined across three areas:
 

\begin{itemize} [leftmargin=0.75cm, noitemsep]
    \item \textbf{Human resources deviation:} To contrast the actual effort with the previous estimated effort of human resources in each task.
    
    \vspace{5pt}
    {\footnotesize
    \centerline{$\displaystyle
    \text{Human resources deviation} =
    \sum_{i \in pit} 
    \left(
    \text{estimated\_cost\_per\_hour}_i - \text{real\_cost\_per\_hour}_i
    \right) 
    \times
    \text{total\_real\_hours}_i
    $}
    }

    {\footnotesize
        Where $i$ indexes each individual in $pit$, the set of all personnel involved in that task.
    }
    
    \item \textbf{Amortization deviation:} To compare the previously estimated use of hardware resources with the actual usage time of those resources.
    
    \vspace{5pt}
    {\footnotesize
    \centerline{$\displaystyle
    \text{Amortization deviation} =
    \sum_{i \in hr} 
    \left(
    \text{estimated\_usage\_hours}_i - \text{real\_usage\_hours}_i
    \right) 
    \times
    \text{price\_per\_hour}_i
    $}
    }
    
    {\footnotesize
        Where $i$ indexes each resource in $hr$, the set of all hardware resources used in the project.
    }

    \item \textbf{Total cost deviation:} To determine each task's real cost deviation from the estimated cost. The general cost described in the formula represents the sum of the personnel and generic costs, excluding contingency and incident costs. Hence, the following formula is applied to each task:

    \vspace{5pt}
    {\small
    \centerline{$\displaystyle
    \text{Total cost deviation} =
    \left(
    \text{estimated\_general\_cost} - \text{real\_general\_cost}
    \right) 
    $}
    }
    
\end{itemize}

\vspace{2pt}

Note that the resulting total cost deviation value can be either positive or negative. A positive value indicates an overestimation of the initial budget for that resource or task. In this case, the additional budget can be reallocated to address issues that arise during the project. Conversely, a negative value indicates an underestimation, requiring contingency and incident costs to cover the budget shortfall.

Furthermore, some costs are not included in the deviation calculation. For example, Azure offers a fixed monthly fee; therefore, overusage will not affect the overall cost deviation. The same applies to the remaining omitted costs. 


